\section{Dual choice and excluded winners}\label{app:choice-duality}

There are three different questions behind replacing union by intersection
in owned choice. Complementing the observation gives an isomorphic algebra
of players excluded from winning. Taking the meet in an owner's existing
choice order gives a different operation, which needs an empty value.
Adjoining the complement dual on the same carrier is a stronger extension:
with the literal formulas on the full powerset, the two families and winner
constants express every operation. We establish these statements before
comparing the game constructions they suggest. Throughout, $K\geq3$.

\subsection{The complement dual}

Write $C(S)=I\setminus S$ and $e_r=I\setminus\{r\}$. Besides the
winner carrier $\Pow=\mathcal P(I)\setminus\{\varnothing\}$, use the
\emph{excluded-winner carrier}
\[
 \mathcal L=\mathcal P(I)\setminus\{I\}.
\]
Complementation is a bijection $C:\Pow\to\mathcal L$, with inverse
given by the same formula. Complementation is not an operation of the winner algebra: every term
operation there is monotone for inclusion (\cref{olv:prop:sat}), and
complement reverses inclusion. De Morgan duality therefore relates two
algebras instead of acting inside one; \cref{dual:thm:complete} shows
what adjoining an internal complement does. For $X,Y\in\mathcal L$, define
\begin{equation}\label{dual:eq:choice}
 X\diamond_r Y=C(C(X)\star_r C(Y))
 =\begin{cases}
 e_r,&X=e_r\text{ or }Y=e_r,\\
 X\cap Y,&\text{otherwise}.
 \end{cases}
\end{equation}
The exception is equality with $e_r$, not merely omission of $r$.
For example, when $I=\{a,b,c\}$, neither $\{b\}$ nor $\{c\}$ is
$e_a$, and $\{b\}\diamond_a\{c\}=\varnothing$.

\begin{proposition}[Transport of the winner algebra]\label{dual:prop:transport}
Each $\diamond_r$ is associative, commutative and idempotent, with
absorbing element $e_r$. Complementation is an isomorphism from the
winner algebra to the algebra on $\mathcal L$ with choices
$\diamond_r$, constants $e_i$, and the usual action of role
permutations. In particular, corresponding signatures have exactly the
same identities, subalgebra and congruence structures, and term
equivalence problems. For every finite nonempty family $(X_j)_{j\in J}$,
\begin{equation}\label{dual:eq:fold}
 \mathop{\diamond_r}_{j\in J}X_j=
 \begin{cases}
 e_r,&X_j=e_r\text{ for some }j,\\
 \bigcap_{j\in J}X_j,&\text{otherwise}.
 \end{cases}
\end{equation}
\end{proposition}
\begin{proof}
An operand's complement is $\{r\}$ precisely when the operand is
$e_r$, and De Morgan's law gives the other case of
\eqref{dual:eq:choice}. Isomorphism transports the semilattice laws,
constants, and identities. A permutation $\pi$ commutes with $C$,
since it is a bijection of $I$. Induction on terms transports their
evaluations, giving the assertions about equivalence and structure.
Complement the finite owned-choice fold to obtain
\eqref{dual:eq:fold}.
\end{proof}

The derived neutral union becomes neutral intersection. Thus for distinct
$r,s,t$ and all $X,Y\in\mathcal L$,
\begin{equation}\label{dual:eq:intersection}
 (X\diamond_r Y)\diamond_t(X\diamond_s Y)=X\cap Y.
\end{equation}
This follows by complementing \cref{game:thm:union}; it is an identity
on the new carrier, not an intersection term on $\Pow$.

\begin{proposition}[Players who never win]\label{dual:prop:semantics}
For a public game $G$, let
\[
 L(G)=\{i:\text{$i$ loses in every SPE of $G$}\}.
\]
Then
\[
 L(G)=C(\den G),\qquad L(\at i)=e_i,\qquad
 L(G\ch r H)=L(G)\diamond_r L(H).
\]
Every proper subset of $I$ occurs as some $L(G)$. Furthermore,
$L(G)=e_r$ exactly when $r$ can force a win in $G$.
\end{proposition}
\begin{proof}
Membership in $C(\den G)$ means that no SPE has winner $i$, which
is exactly the displayed universal condition. Backward induction and
\eqref{dual:eq:choice} give composition. Every nonempty winner set is
realizable by \cref{game:prop:winner-budget}; complementing proves
realizability here. The last assertion is the winner-atom criterion.
\end{proof}

If a branch has excluded-winner value $e_r$, its only possible SPE
winner is $r$, so the owner ensures that every other player loses.
Otherwise a player is excluded at the parent exactly when it is excluded
in both branches. This is the intersection clause. The empty value says
that every player wins in some SPE; the excluded value $I$ would say
that there is no SPE. For example, an $a$-choice between terminal wins
for $b$ and $c$ has excluded-winner value
\[
 \{a,c\}\diamond_a\{a,b\}=\{a\}
 \qquad(I=\{a,b,c\}).
\]

These are universal losses, not possible losses. The set of players who
lose in \emph{some} SPE is $e_i$ when $\den G=\{i\}$ and is $I$
when $\den G$ is ambiguous. That observation discards the distinction
between ambiguous values, whereas $L$ retains the entire winner value.

\subsection{Order duality and the actual meet}

The choice order on $\Pow$ is
\[
 S\sqsubseteq_r T
 \quad\Longleftrightarrow\quad
 T=\{r\}\ \text{or}\ (S\subseteq T\text{ and }S\ne\{r\}).
\]
The dual choice is meet for a different order on $\mathcal L$:
\begin{equation}\label{dual:eq:order}
 X\preceq_r Y
 \quad\Longleftrightarrow\quad
 X=e_r\ \text{or}\ (X\subseteq Y\text{ and }Y\ne e_r).
\end{equation}
Indeed $X\preceq_r Y$ is equivalent to
$C(Y)\sqsubseteq_r C(X)$. Thus $e_r$ has been moved to the bottom
of the inclusion order, just as $\{r\}$ was moved to the top.
The original order on nonempty sets has no meet for two distinct
nonowner atoms: they have no lower bound.

For a comparison on a common carrier, extend both literal formulas to
all of $\mathcal P(I)$. Denote these operations by $\star_r^+$ and
$\diamond_r^+$: their definitions are respectively
\eqref{game:eq:own} and \eqref{dual:eq:choice}, now allowing every
subset. They remain complement conjugates. Here
$\varnothing$ is an identity for $\star_r^+$ and $I$ is an identity
for $\diamond_r^+$; their absorbers remain $\{r\}$ and $e_r$.

\begin{proposition}[Meet in the choice order]\label{dual:prop:meet}
On the full powerset, the order $\sqsubseteq_r$ is a bounded lattice,
with bottom $\varnothing$, top $\{r\}$, join $\star_r^+$, and meet
\begin{equation}\label{dual:eq:actual-meet}
 M_r(S,T)=\begin{cases}
 T,&S=\{r\},\\
 S,&T=\{r\},\\
 \varnothing,&S,T\ne\{r\}\text{ and }S\cap T=\{r\},\\
 S\cap T,&\text{otherwise}.
 \end{cases}
\end{equation}
This meet differs from $\diamond_r^+$. The pair
$(\star_r^+,\diamond_r^+)$ does not satisfy the lattice absorption
laws.
\end{proposition}
\begin{proof}
When one input is top, the other is the greatest lower bound. Otherwise
the lower bounds are exactly the subsets of $S\cap T$ other than
$\{r\}$. If the intersection is $\{r\}$, only $\varnothing$
remains; in every other case the intersection itself is greatest.
The join assertion follows from the definition of $\sqsubseteq_r$.
For distinct $r,b,c$, take $S=\{b\}$ and $T=\{r\}$.
Then $M_r(S,T)=S$, but
\[
 S\diamond_r^+(S\star_r^+T)
 =\{b\}\diamond_r^+\{r\}=\varnothing\ne S.
\]
Neither singleton is $e_r$, since $K\geq3$.
\end{proof}

Replacing union by intersection while keeping $\{r\}$ absorbing
does not give a semilattice even after allowing the empty set. If
$S\mathbin{\otimes_r}T$ is $\{r\}$ when an input is $\{r\}$
and $S\cap T$ otherwise, take distinct $r,b,c$ and put
$S=\{r,b\}$, $T=\{r,c\}$, $U=\{b\}$. Then
\[
 (S\mathbin{\otimes_r}T)\mathbin{\otimes_r}U=\{r\},\qquad
 S\mathbin{\otimes_r}(T\mathbin{\otimes_r}U)=\varnothing.
\]
An ordinary intersection can create the exceptional singleton from two
nonexceptional inputs, which then changes the next stage. The coatom
exception in the complement dual avoids precisely this failure: an
intersection of proper sets different from $e_r$ cannot equal $e_r$.

The full-powerset extension just used is algebraic. It differs from the
annihilating-zero extension in \cref{query:prop:relations} and from
equilibrium composition with an equilibrium-free continuation in
\cref{olv:prop:infinite-grafting}. In the latter situation, every
off-path continuation must also carry an SPE, so an empty child makes
the parent empty. In contrast,
$\varnothing\star_r^+T=T$. Complementing the annihilating-zero
extension gives an operation with $I$ absorbing; it agrees with
$\diamond_r$ on $\mathcal L$ but differs from
$\diamond_r^+$ at $I$. Thus neither typed algebra determines how the
other endpoint should behave when the carriers are identified.

Even identities of the original reduct can fail at the new endpoint.
For $b\ne r$, distribution over union fails because
\[
 \{b\}\star_r^+(\varnothing\cup\{r\})=\{r\},\qquad
 (\{b\}\star_r^+\varnothing)\cup
 (\{b\}\star_r^+\{r\})=\{b,r\}.
\]
On nonempty inputs that distributive identity holds; complementing it
gives distribution of $\diamond_r$ over intersection on $\mathcal L$.

\subsection{Adjoining both literal operations}

Fix the literal extensions from the preceding subsection. Let
\[
 \mathbf D_I=(\mathcal P(I);
       (\star_r^+,\diamond_r^+)_{r\in I},(\{i\})_{i\in I}).
\]
There is no role-renaming operation in this signature. The neutral-union
term and its complement dual define ordinary union and intersection on
\emph{all} subsets. To see that neutral union still works at the new
endpoint, if one argument is empty then both inner choices return the
other argument, and idempotence finishes the calculation. The other
cases are \cref{game:thm:union}. Complementation gives intersection,
including at $I$. Consequently $\varnothing=\{a\}\cap\{b\}$
for $a\ne b$, and $I=\bigcup_{i\in I}\{i\}$ are term constants.

\begin{lemma}[Transfer of Boolean tests]\label{dual:lem:transfer}
For distinct roles $p,q$, the term
\begin{equation}\label{dual:eq:transfer}
 N_{pq}(U)=(U\star_p^+\{q\})\cap\{q\}
\end{equation}
satisfies $N_{pq}(\varnothing)=\{q\}$ and
$N_{pq}(\{p\})=\varnothing$.
For any roles $p,q$, there are terms $P_{pq}$ and $Q_{pq}$ on
$\{\varnothing,\{p\}\}$ such that
\[
 P_{pq}(U)=\begin{cases}\{q\},&U=\{p\},\\
 \varnothing,&U=\varnothing,\end{cases}
 \qquad
 Q_{pq}(U)=\begin{cases}\varnothing,&U=\{p\},\\
 \{q\},&U=\varnothing.\end{cases}
\]
\end{lemma}
\begin{proof}
The two evaluations of $N_{pq}$ follow from the identity and absorber
laws for $\star_p^+$. If $p\ne q$, choose $h$ different from
both and put $P_{pq}=N_{hq}\circ N_{ph}$; if $p=q$, use the
identity. For $p\ne q$, put $Q_{pq}=N_{pq}$. For $p=q$, choose
distinct $h,k$ outside $\{p\}$ and put
$Q_{pp}=N_{kp}\circ N_{hk}\circ N_{ph}$.
Two transfers preserve the truth bit and three negate it.
\end{proof}

In particular, complement itself is a term of $\mathbf D_I$:
\begin{equation}\label{dual:eq:term-complement}
 C(S)=\bigcup_{r\in I}Q_{rr}(S\cap\{r\}).
\end{equation}
Unlike the typed complementation $\Pow\leftrightarrow\mathcal L$,
this is an internal operation of the enlarged algebra.

\begin{theorem}[Functional completeness]\label{dual:thm:complete}
For every finite $I$ with $|I|\geq3$ and every finite arity $n$,
every function
\[
 f:\mathcal P(I)^n\longrightarrow\mathcal P(I)
\]
is a term operation of $\mathbf D_I$.
\end{theorem}
\begin{proof}
Fix a marker $a\in I$. The term
$P_{ra}(S_j\cap\{r\})$ is $\{a\}$ precisely when
$r\in S_j$, and $Q_{ra}(S_j\cap\{r\})$ is $\{a\}$
precisely when $r\notin S_j$. For a fixed input tuple
$A=(A_1,\ldots,A_n)$, intersect the positive or negative test for
each pair $(j,r)$ according to whether $r\in A_j$. Starting the
intersection at $\{a\}$ gives a term $E_A$ with
\[
 E_A(S)=\begin{cases}\{a\},&S=A,\\
 \varnothing,&S\ne A.\end{cases}
\]
This also covers $n=0$, when there is one input tuple and no tests.
Now take the finite union
\begin{equation}\label{dual:eq:truth-table}
 \bigcup_{A\in\mathcal P(I)^n}\,
 \bigcup_{r\in f(A)} P_{ar}(E_A(S)),
\end{equation}
using the empty-set constant for an empty union. Exactly the summand
$A=S$ contributes, and its contribution is $f(S)$.
All intersections, unions, tests, and constants used here are terms
in the specified signature.
\end{proof}

The role threshold is sharp. With one or two roles the singleton values
form a subalgebra for both literal choice families. For two roles,
$e_r$ is the other singleton, so $\diamond_r^+$ on singleton inputs
is the other owner's original choice; with one role, idempotence
suffices. Thus a term with singleton constants sends every all-singleton
input tuple to a singleton and cannot express the constant empty
function. The proof for $K\geq3$ also shows that ordinary intersection,
adjoined to owned choices and singleton constants, already suffices for
functional completeness: dual choices are used only to define it.

Thus this extension removes the restrictions on expressible operations
in the original algebra. For example it has the discriminator
$d(X,Y,Z)=Z$ if $X=Y$ and $d(X,Y,Z)=X$ otherwise, and the ternary
majority operation
$(X\cap Y)\cup(Y\cap Z)\cup(Z\cap X)$.
The latter is a near-unanimity operation, in contrast with
\cref{gen:cor:trace}. Functional completeness also implies that
$\mathbf D_I$ has only equality and the universal congruence: if a
congruence identifies distinct $X,Y$, an arbitrary unary term can
send $X,Y$ to any prescribed pair of values.

The endpoint convention is essential to the proof: the negating transfer
uses $\varnothing\star_p^+\{q\}=\{q\}$.
The theorem is about the specified algebra $\mathbf D_I$, rather than
a representation of arbitrary functions by the original public games.

\subsection{Alternative game readings}

\paragraph{Reversing an owner's preference.}
Changing a player's objective changes which branch is optimal; it does
not change existential enumeration of equilibrium outcomes into
intersection. More generally, allow each player $i$ a fixed set $A_i$
of preferred terminal labels, with payoff one inside $A_i$ and zero
outside. For an $i$-node write $A=A_i$. The original game is the case
$A_i=\{i\}$.

Possible SPE winners then compose by the shared-objective choice
$\star_A$ of \cref{context:prop:collective}: the union of the branches,
filtered by $A$ when some branch lies inside $A$.

If the owner prefers its own loss, take $A=I\setminus\{r\}$.
Then the output is $(S\cup T)\setminus\{r\}$ when $r$ is absent
from either child, and is $S\cup T$ otherwise. In particular, a
choice between wins for $b$ and $c$, both different from $r$, still
has outcome set $\{b,c\}$.

Complementing \eqref{context:eq:collective} gives a useful explanation of
the exceptional coatom in \eqref{dual:eq:choice}. If $B=I\setminus A$,
then the excluded-outcome operation on proper subsets is
\begin{equation}\label{dual:eq:objective-dual}
 (X,Y)\longmapsto
 \begin{cases}
 (X\cap Y)\cup B,&B\subseteq X\text{ or }B\subseteq Y,\\
 X\cap Y,&\text{otherwise}.
 \end{cases}
\end{equation}
For $A=\{r\}$, $B=e_r$ and a proper set contains $B$ precisely
when it equals $B$. In the exceptional case $(X\cap Y)\cup B=B$,
recovering exactly $\diamond_r$.

\paragraph{Requiring a common winner.}
Intersection has a direct interpretation as a compatibility condition on
\emph{two} equilibrium witnesses. Define
\[
 \operatorname{Sync}(G,H)=
 \{i:\exists\sigma\in\operatorname{SPE}(G),\
            \exists\tau\in\operatorname{SPE}(H),\
            \operatorname{out}(\sigma)=\operatorname{out}(\tau)=i\}.
\]
By independent choice of the two witnesses,
\begin{equation}\label{dual:eq:sync}
 \operatorname{Sync}(G,H)=\den G\cap\den H.
\end{equation}
This relation can be empty even for two terminal games. It records
compatible equilibrium witnesses; a game implementing agreement would
also need actions and payoffs for disagreement. Those choices can change
which profiles are equilibria.

There is no ordinary public-game term implementing even the nonempty
cases of \eqref{dual:eq:sync}, by
\cref{query:prop:no-intersection}. Similarly the complement dual on
possible-winner sets is not a public-game constructor: for distinct
$r,b,c$, the nonempty singleton inputs $\{r\},\{b\}$ already
give $\varnothing$ under $\diamond_r$. In the excluded-winner
reading this causes no failure of equilibrium existence: the same
empty value means that no player is excluded. The game interpretation
therefore depends on retaining the distinction between winner sets,
excluded-winner sets, and compatibility constraints.
